\chapter{全纯函数}\label{ch:b3:holomorphic}

复可微性看起来只是实理论上的一个小变化，
无非一个极限、一个商。它实则是
另一个世界。在某开集上\emph{一次}可微的函数
（$\C$ 的开子集）自动无穷次可微，
解析，由其在某一
点附近的值唯一确定于整个区域，并受刚性的全局原理
约束（刘维尔、最大模原理）。这一切源于一个奇迹，
即柯西定理：全纯函数
沿星形区域中的闭路径积分为零。
本章证明这一奇迹（采用古尔萨的论证，
除可微性外不假设任何额外的正则性），并收获它的
经典推论，并证明本书
自\cref{ch:b3:galois}起一直借用的定理：每个非常值
复多项式都有根。

下文约定 $\Omega \subseteq \C$ 为开集，$D(a, r)$
表示开圆盘。

\section{复可微性}

\begin{definition}\label{def:b3:holomorphic:holo}
$f \colon \Omega \to \C$ 称为\emph{全纯函数}\index{全纯
函数}在 $\Omega$ 上，若对每个 $z_0 \in \Omega$
\[
f'(z_0) = \lim_{h\to0}\frac{f(z_0 + h) - f(z_0)}{h}
\]
极限存在（$h \in \C^*$）。全纯函数的和、积、
商（分母非零）、复合亦为全纯，
并满足通常的求导公式（本科一、二年级的证明
原样照搬，它们只用到域的运算与极限）。
$\mathcal H(\Omega)$ 表示
$\Omega$ 上全纯函数的集合。
\end{definition}

\begin{proposition}[柯西–黎曼]
\label{prop:b3:holomorphic:cauchyriemann}
记 $f(x + \iu y) = P(x,y) + \iu Q(x,y)$，则 $f$
在 $z_0$ 处全纯，当且仅当 $f$（作为两个实变量的映射）
$\R$-可微于 $z_0$，且满足\index{柯西–黎曼
方程}：
\[
\frac{\partial P}{\partial x} = \frac{\partial Q}{\partial y},
\qquad
\frac{\partial P}{\partial y} = -\frac{\partial Q}{\partial x}
\qquad \text{在 } z_0 ;
\]
等价地，实微分是乘以复数
$f'(z_0)$ 的运算。
\end{proposition}

\begin{proof}
$\C$-可微性说明 $f(z_0 + h) = f(z_0) + ch +
o(\abs h)$，其中 $c = f'(z_0)$：一个 $\R$-线性微分
它即乘以 $c = a + \iu b$ 的运算，其矩阵
为 $\bigl(\begin{smallmatrix} a & -b\\ b & a
\end{smallmatrix}\bigr)$（基为 $(1, \iu)$），恰好
即偏导数所满足的上述关系。反之，这样的
微分是 $\C$-线性的，且 $o(\abs h)$ 的定义
相互吻合。
\end{proof}

\begin{example}\label{ex:b3:holomorphic:examples}
关于 $z$ 的多项式，极点之外的有理函数，以及
——由本科二年级关于幂级数的逐项求导定理（其证明
在 $\C$ 上逐字不变）得知，每个
幂级数 $\sum a_n(z - a)^n$ 在其收敛圆盘内的和
均为全纯函数，导数为 $\sum na_n(z -
a)^{n-1}$（收敛半径不变）。特别地，$\exp z = \sum z^n/n!$
是整函数（$\C$ 上的全纯函数），且 $\exp' = \exp$。
另一方面，$z \mapsto \bar z$、$\abs z$、$\operatorname{Re}z$
处处不全纯（柯西–黎曼方程处处不成立）：
全纯性是保定向、保角的刚性，而非
光滑性。
\end{example}

\section{围道积分}

\begin{definition}\label{def:b3:holomorphic:contour}
称\emph{路径}为分段 $\mathcal C^1$ 映射 $\gamma \colon
\intcc ab \to \C$；若 $\gamma(a) =
\gamma(b)$，则称\emph{闭路径}。设
$f$ 为 $\gamma$ 的像上的连续函数：\index{围道积分}
\[
\int_\gamma f(z)\,\dd z =
\int_a^b f(\gamma(t))\,\gamma'(t)\,\dd t,
\qquad
\Bigl|\int_\gamma f\,\dd z\Bigr| \leq
\sup_{\gamma}\abs f\cdot\operatorname{length}(\gamma)
\]
（\emph{ML 不等式}；长度 $= \int_a^b\abs{\gamma'}$）。
积分对保持定向的 $\mathcal C^1$
重参数化不变，且在定向反转时变号。
\end{definition}

\begin{proposition}[原函数]\label{prop:b3:holomorphic:primitive}
设 $f$ 为 $\Omega$ 上的连续函数，下列条件等价：
(i) $f$ 有\emph{原函数} $F \in \mathcal H(\Omega)$（$F'
= f$）；(ii) $\int_\gamma f\,\dd z = 0$ 对每条闭路径
$\gamma$ 在 $\Omega$ 中成立。在此情形下 $\int_\gamma f\,\dd z =
F(\gamma(b)) - F(\gamma(a))$ 对每条路径成立。
\end{proposition}

\begin{proof}
(i)$\Rightarrow$：$\frac{\dd}{\dd t}F(\gamma(t)) =
F'(\gamma(t))\gamma'(t)$（链式法则，分段成立），故
积分相消为端点之差；闭路径
给出 $0$。(ii)$\Rightarrow$(i)：固定某一
连通分支中的一点 $z_*$，定义 $F(z) = \int_{\gamma_z}f\,\dd z$，沿任一
从 $z_*$ 到 $z$ 的路径（良定义：两条路径之差为
一条闭路径）；对充分小的 $h$，取从 $z$ 到
$z + h$ 的线段，
\[
\frac{F(z + h) - F(z)}{h} - f(z)
= \frac1h\int_{[z, z+h]}\bigl(f(w) - f(z)\bigr)\dd w
\longrightarrow 0
\]
由 ML 不等式及 $f$ 在 $z$ 处的连续性。
\end{proof}

\begin{definition}[环绕数]\label{def:b3:holomorphic:index}
对闭路径 $\gamma$ 及 $z \notin
\operatorname{im}\gamma$，其\emph{环绕数}\index{环绕数}
\index{路径的环绕数}为
\[
\operatorname{Ind}_\gamma(z) = \frac1{2\iu\pi}
\int_\gamma\frac{\dd w}{w - z} .
\]
它是\emph{整数}：令 $\varphi(t) =
\int_a^t\frac{\gamma'(s)}{\gamma(s) - z}\dd s$，函数
$(\gamma(t) - z)\eu^{-\varphi(t)}$ 的导数为零，
（分段），故为常数；取 $t = b$，$\eu^{\varphi(b)} =
\frac{\gamma(b) - z}{\gamma(a) - z} = 1$，故 $\varphi(b) \in
2\iu\pi\Z$。作为 $z$ 的函数，环绕数在
$\C\setminus\operatorname{im}\gamma$ 上连续（由控制收敛定理），
从而在每个连通分支上为常数，且在无界分支上为 $0$（由
ML 不等式，积分趋于 $0$，当 $z \to
\infty$ 时）。对圆 $\gamma(t) = a + r\eu^{\iu t}$，$t
\in \intcc0{2\pi}$ 时，$\operatorname{Ind}_\gamma(z) = 1$，对 $z
\in D(a,r)$（在 $z = a$ 处计算：$\frac1{2\iu\pi}
\int_0^{2\pi}\frac{r\iu\eu^{\iu t}}{r\eu^{\iu t}}\dd t = 1$；
由常数性即可得其余部分）。
\end{definition}

\section{柯西定理}

\begin{theorem}[古尔萨定理]\label{thm:b3:holomorphic:goursat}
设 $f \in \mathcal H(\Omega)$，$T \subseteq \Omega$ 为
闭三角形区域。则 $\int_{\partial T}f\,\dd z = 0$
（边界绕行一次，方向任意）。\index{古尔萨定理}
\end{theorem}

\begin{proof}
设 $I(T) = \int_{\partial T}f\,\dd z$。连接
各边中点，将 $T$ 分为四个边长减半的三角形 $T^{(1)},
\dots, T^{(4)}$，内部的边两两抵消：$I(T) =
\sum_iI(T^{(i)})$。在其中选取 $T_1$，使
$\abs{I(T_1)} \geq \frac14\abs{I(T)}$，并迭代此过程，得一嵌套
序列 $T \supseteq T_1 \supseteq T_2 \supseteq\cdots$，满足
\[
\abs{I(T_n)} \geq 4^{-n}\abs{I(T)},
\qquad
\operatorname{diam}T_n = 2^{-n}\operatorname{diam}T,
\quad
\operatorname{length}(\partial T_n) =
2^{-n}\operatorname{length}(\partial T).
\]
交集 $\bigcap T_n$ 为单点 $z_0$
（嵌套紧集，直径趋于零，
由\cref{thm:b3:topology:compactprops}(3)）。在
$z_0$ 处可导：任给 $\varepsilon$，当 $n$ 充分大时，在 $T_n$ 上，
\[
f(z) = f(z_0) + f'(z_0)(z - z_0) + R(z),
\qquad \abs{R(z)} \leq \varepsilon\abs{z - z_0} \leq
\varepsilon\operatorname{diam}T_n .
\]
仿射部分有原函数，其在闭
$\partial T_n$ 上的积分为零（见\cref{prop:b3:holomorphic:primitive}），
于是余下
\[
\abs{I(T_n)} = \Bigl|\int_{\partial T_n}R\Bigr|
\leq \varepsilon\operatorname{diam}(T_n)\,
\operatorname{length}(\partial T_n)
= \varepsilon\,4^{-n}\operatorname{diam}(T)
\operatorname{length}(\partial T) .
\]
与 $\abs{I(T_n)} \geq 4^{-n}\abs{I(T)}$ 比较：
$\abs{I(T)} \leq \varepsilon\cdot\text{常数}$ 对每个
$\varepsilon$，故 $I(T) = 0$。
\end{proof}

\begin{theorem}[柯西定理，星形区域版本]
\label{thm:b3:holomorphic:cauchy}
设 $\Omega$ 关于 $c$ 是\emph{星形的}（即每条线段
$[c, z]$，$z \in \Omega$，均含于 $\Omega$） --- 例如凸集。
每个 $f \in \mathcal H(\Omega)$ 在 $\Omega$ 上均有原函数；
从而 $\int_\gamma f\,\dd z = 0$ 对\emph{每条}
闭路径 $\gamma$ 在 $\Omega$ 中成立。\index{柯西定理}
\end{theorem}

\begin{proof}
定义 $F(z) = \int_{[c,z]}f\,\dd w$。对 $z, z + h \in
\Omega$ 且 $[z, z+h] \subseteq \Omega$（对充分小的 $h$ 成立），
则以 $c, z, z+h$ 为顶点的三角形含于 $\Omega$
（星形性：其每一点都落在某线段 $[c,
w]$，其中 $w \in [z, z+h] \subseteq \Omega$）。由古尔萨定理得
\[
F(z + h) - F(z) = \int_{[z, z+h]}f\,\dd w,
\]
并由差商计算
（见\cref{prop:b3:holomorphic:primitive}）得 $F' = f$。闭
路径积分为零同样由同一
命题推出。
\end{proof}

\begin{theorem}[柯西积分公式]
\label{thm:b3:holomorphic:formula}
设 $f \in \mathcal H(\Omega)$，$\bar D(a, r) \subseteq
\Omega$，$C_r$ 为圆周 $\partial D(a,r)$，沿
逆时针方向绕行一次。则对每个 $z \in D(a, r)$：\index{柯西
积分公式}
\[
f(z) = \frac1{2\iu\pi}\int_{C_r}\frac{f(w)}{w - z}\,\dd w .
\]
\end{theorem}

\begin{proof}
固定 $z$，在 $\Omega$ 上定义
\[
g(w) = \begin{cases}
\dfrac{f(w) - f(z)}{w - z} & w \neq z,\\[4pt]
f'(z) & w = z :
\end{cases}
\]
$g$ 在 $\Omega$ 上连续，且除 $z$ 外全纯。古尔萨定理
对 $g$ 在每个三角形 $T \subseteq \Omega'$ 上成立，其中
$\Omega'$ 为略大于 $\bar D(a,r)$ 而位于
$\Omega$ 内的圆盘：若 $z \notin T$，直接成立；若 $z \in T$，则将 $T$
分成以 $z$ 为顶点的小三角形及
不含 $z$ 的三角形；在以 $z$ 为顶点的三角形上，ML 不等式给出
$\abs{\int_{\partial T'}g} \leq \sup_{T'}\abs
g\cdot\operatorname{length} \to 0$，三角形缩小时，
其余部分由古尔萨定理知积分为零 --- 因此
$\int_{\partial T}g = 0$ 在一切情形下成立。
\cref{thm:b3:holomorphic:cauchy} 的证明只用到了此三角形
性质：$g$ 在凸集 $\Omega'$ 上有原函数，因此
$\int_{C_r}g = 0$，即
\[
\frac1{2\iu\pi}\int_{C_r}\frac{f(w)}{w - z}\dd w
= f(z)\,\frac{1}{2\iu\pi}\int_{C_r}\frac{\dd w}{w - z}
= f(z)\operatorname{Ind}_{C_r}(z) = f(z) .
\]
\end{proof}

\section{解析性及其连锁}

\begin{theorem}[全纯 $=$ 解析]
\label{thm:b3:holomorphic:analytic}
设 $f \in \mathcal H(\Omega)$，且 $D(a, R) \subseteq \Omega$。
则\index{全纯函数的解析性}
\[
f(z) = \sum_{n\geq0}c_n\,(z - a)^n \quad \text{在 } D(a, R),
\qquad
c_n = \frac1{2\iu\pi}\int_{C_r}\frac{f(w)}{(w -
a)^{n+1}}\,\dd w \ \ (0 < r < R),
\]
这些系数与 $r$ 无关。因此 $f$ 是
无穷次 $\C$-可微，$c_n = f^{(n)}(a)/n!$，并且
\emph{柯西估计}\index{柯西估计}成立：
\[
\abs{c_n} \;\leq\; \frac{\sup_{\abs{w - a} =
r}\abs{f(w)}}{r^{n}} .
\]
\end{theorem}

\begin{proof}
当 $\abs{z - a} < r$ 时，将柯西核展开为
几何级数：
\[
\frac1{w - z} = \frac1{(w - a)\bigl(1 - \frac{z - a}{w -
a}\bigr)} = \sum_{n\geq0}\frac{(z - a)^n}{(w - a)^{n+1}},
\]
该级数关于 $w$ 在 $C_r$ 上正规收敛（$\abs{\frac{z-a}{w-a}} =
\frac{\abs{z-a}}r < 1$）：将级数逐项
乘以 $\frac{f(w)}{2\iu\pi}$ 后积分（一致
收敛保证了交换次序的合理性），并应用\cref{thm:b3:holomorphic:formula}。
幂级数无穷次可微，且 $c_n =
f^{(n)}(a)/n!$（本科二年级），这也说明 $c_n$
与 $r$ 无关。估计式：对系数积分
用 ML 估计即可。
\end{proof}

\begin{example}[奇点决定收敛半径]
\label{ex:b3:holomorphic:radius}
为什么看似无害的实函数 $\frac1{1 + x^2}$ 在 $x = 3$ 处的泰勒级数
仅当 $\abs{x - 3} <
\sqrt{10}$ 时才收敛，而实轴上并无异常？原因在于
上述定理使 $a$ 处的收敛半径
等于 $a$ 到最近的全纯性失效点的距离。
这里 $f(z) = \frac1{1 + z^2}$ 在
$\C\setminus\{\pm\iu\}$ 上恰好全纯，故
$a = 3$ 处的展开在最大的圆盘上收敛，该圆盘避开
$\pm\iu$，半径为 $\abs{3 - \iu} = \sqrt{10}$——而
不可能在更大的圆盘上收敛，因为该和会将
$f$ 全纯地延拓到 $\pm\iu$ 的邻域上，而在那里
$\abs f \to \infty$。实理论只看到神秘的
半径 $\sqrt{10}$；复平面则看到两个极点。
这就是实用法则：\emph{求收敛半径，
只需找出奇点}——例如
$\tan$ 在 $0$ 处的泰勒级数的收敛半径为 $\frac\pi2$（最近的
$\cos$ 的零点），而伯努利母函数
$\frac z{\eu^z - 1}$（\cref{pb:b3:holomorphic:1}，第 VI 部分）
的收敛半径为 $2\pi$（$\eu^z - 1$ 的最近非零零点：
$\pm2\iu\pi$）。
\end{example}

\begin{corollary}[刘维尔；达朗贝尔--高斯]
\label{cor:b3:holomorphic:liouville}
有界整函数必为常数。因此
$\C$ 上每个非常数多项式都有根，即 \emph{$\C$ 是
代数闭的}。\index{刘维尔定理}\index{代数基本
定理}
\end{corollary}

\begin{proof}
若 $\abs f \leq M$ 在 $\C$ 上：对每个 $a$ 与 $r$，$\abs{c_1(a)}
= \abs{f'(a)} \leq M/r \to 0$：于是 $f' \equiv 0$，从而 $f$
为常数（在连通的 $\C$ 上：导数为零蕴含
局部常数——沿线段积分）。若 $P$ 没有
根，则 $1/P$ 为整函数且有界（$\abs{P(z)} \to
\infty$（$\abs z \to \infty$ 时）：首项占主导，所以
$\abs{1/P}$ 在大圆盘之外很小，并在紧的圆盘上连续）：
为常数，与 $P$ 非常数矛盾。
（周末问题给出第二个初等证明，以及
其代数推论。）
\end{proof}

\begin{theorem}[零点孤立；恒等定理]
\label{thm:b3:holomorphic:identity}
设 $\Omega$ 为\emph{连通}的，$f \in \mathcal
H(\Omega)$，$f \not\equiv 0$。则每个零点 $a$（指 $f$ 的零点）
具有有限的\emph{阶}：$f(z) = (z - a)^m\,g(z)$，其中 $g \in
\mathcal H(\Omega)$，$g(a) \neq 0$，且 $f$ 的零点
在 $\Omega$ 中无聚点。因此，若
定义在 $\Omega$ 上的两个全纯函数在某个集合上相等，
该集合在 $\Omega$ 中有聚点，则它们
处处相等。\index{恒等定理}\index{孤立零点}
\end{theorem}

\begin{proof}
设 $Z$ 是使 $f$ 的所有导数
都为零的点的集合。$Z$ 是闭的（闭集之交），且是开的：
若所有 $c_n = 0$ 于 $a$ 处，则幂级数展开使 $f
\equiv 0$ 在 $a$ 附近的一个圆盘上。由连通性，$Z =
\varnothing$ 或 $Z = \Omega$，而后者被 $f
\not\equiv 0$ 排除。因此在零点 $a$ 处有某个系数非零：
取 $m$ 为使 $c_m \neq 0$ 的最小值，则 $f(z) = (z -
a)^m\sum_{n\geq m}c_n(z-a)^{n-m}$ 在某个圆盘上成立，而该级数
定义了 $g$，它在 $a$ 附近全纯，且 $g(a) = c_m \neq 0$；
令 $g = f/(z-a)^m$ 给出 $a$ 以外的延拓（该处全纯）。由于
$g(a) \neq 0$ 且 $g$ 连续，所以 $f$ 在
$a$ 的一个邻域内没有其他零点：零点是孤立的，而孤立点所成的
集合在 $\Omega$ 中无聚点（零点的聚点
必为零点（由连续性），且不是孤立的）。
恒等性：对两者之差应用，
其零点集有聚点，从而迫使它落入
$Z = \Omega$ 这一分支。
\end{proof}

\begin{theorem}[均值与最大模]
\label{thm:b3:holomorphic:maximum}
设 $f \in \mathcal H(\Omega)$。
\begin{enumerate}
  \item （均值）若 $\bar D(a, r) \subseteq \Omega$：
        $f(a) = \frac1{2\pi}\int_0^{2\pi}f(a +
        r\eu^{\iu t})\,\dd t$.
  \item （最大模原理）若 $\Omega$ 连通，且
        $\abs f$ 在某点取得局部极大值，
        即该点位于 $\Omega$ 内，则 $f$ 为常函数。由此，若
        $\Omega$ 有界，且 $f$ 在
        $\bar\Omega$ 上连续，则 $\sup_{\bar\Omega}\abs f =
        \sup_{\partial\Omega}\abs f$。\index{最大模原理}
\end{enumerate}
\end{theorem}

\begin{proof}
(1) 即圆心处的柯西公式：将 $C_r$ 参数化。
(2) 设 $\abs f \leq \abs{f(a)}$ 在 $\bar D(a, \rho)
\subseteq \Omega$ 上成立。若 $f(a) = 0$，则 $f \equiv 0$ 于 $a$ 附近。
否则，对 $0 < r \leq \rho$，由均值公式得
\[
\abs{f(a)} \leq \frac1{2\pi}\int_0^{2\pi}\abs{f(a +
r\eu^{\iu t})}\,\dd t \leq \abs{f(a)} :
\]
连续非负的被积函数 $\abs{f(a)} - \abs{f(a +
r\eu^{\iu t})}$ 的平均值为零，从而恒为零，于是 $\abs f$ 在圆盘上
为\emph{常数} $= \abs{f(a)} \ne 0$。
圆盘上模为非零常数的全纯函数必为常数：
对 $P^2 + Q^2 = \text{常数}$ 求导，得 $PP_x + QQ_x
= 0$，且 $PP_y + QQ_y = 0$；将柯西–黎曼关系
$P_y = -Q_x$，$Q_y = P_x$ 代入第二式，
得线性方程组：
\[
P\,P_x + Q\,Q_x = 0, \qquad -P\,Q_x + Q\,P_x = 0,
\]
其行列式为 $P^2 + Q^2 \neq 0$，于是 $P_x = Q_x = 0$，
$f' = P_x + \iu Q_x = 0$，故 $f$ 在圆盘上为常数。
由恒等定理，常数性推广至整个 $\Omega$。边界部分：$\abs f$
在紧集 $\bar\Omega$ 上取得上确界；若极大值在内部取得，
则 $f$ 为常数，因此上确界在
边界上取得（任何情形）。
\end{proof}

\begin{theorem}[魏尔斯特拉斯收敛定理]
\label{thm:b3:holomorphic:weierstrassconv}
若 $f_n \in \mathcal H(\Omega)$ 一致收敛于 $f$，
在 $\Omega$ 的每个紧子集上成立，则 $f \in \mathcal
H(\Omega)$，且 $f_n^{(k)} \to f^{(k)}$ 在紧集上一致收敛，
对每个 $k$ 成立。\index{魏尔斯特拉斯收敛定理}
\end{theorem}

\begin{proof}
$f$ 连续。对任一闭三角形 $T \subseteq
\Omega$: $\int_{\partial T}f = \lim\int_{\partial T}f_n = 0$
（在紧集 $\partial T$ 上一致收敛；对
$f_n$ 用古尔萨定理）。由\cref{thm:b3:holomorphic:cauchy}的论证，
$f$ 有局部原函数（圆盘是凸的，仅用到了三角形性质），
即局部地有 $f = F'$，其中 $F$
全纯，且 $F$ 解析
（\cref{thm:b3:holomorphic:analytic}），从而 $f = F'$ 亦
全纯。（此即
\emph{莫雷拉定理}\index{莫雷拉定理}：连续且三角形积分为零者
必全纯。）关于导数，对
$\bar D(a, 2r) \subseteq \Omega$ 及 $z \in \bar D(a, r)$，
由导数的柯西公式（在积分号下对
\cref{thm:b3:holomorphic:formula}求导，或用
系数公式）得
\[
\abs{f_n'(z) - f'(z)} =
\Bigl|\frac{1}{2\iu\pi}\int_{C_{2r}}\frac{f_n(w) - f(w)}{(w -
z)^2}\,\dd w\Bigr|
\leq \frac{2r\,\sup_{C_{2r}}\abs{f_n - f}}{r^2} \to 0
\]
在 $\bar D(a,r)$ 上一致成立；以有限个此类圆盘覆盖紧集，
并对更高的 $k$ 反复应用。
\end{proof}

\begin{omfigure}
\begin{tikzpicture}[scale=1.5]
  \draw[->] (-1.7,0) -- (1.9,0);
  \draw[->] (0,-1.35) -- (0,1.5);
  \draw[omDef, very thick] (0,0) circle (1);
  \foreach \a in {0, 30, 60, 90, 120, 150}
    \draw[omDef!50] (\a:1.06) -- (\a:0.94)
      (\a+180:1.06) -- (\a+180:0.94);
  \fill[omThm] (0.35,0.25) circle (1pt)
    node[above, font=\small] {$z$};
  \fill (0,0) circle (1pt) node[below left, font=\small] {$a$};
  \draw[omProp, thick, ->]
    (1,0) arc (0:75:1);
  \node[omProp, font=\small] at (60:1.28) {$C_r$};
  \node[font=\small, omThm] at (1.45,0.75)
    {$f(z) = \dfrac1{2\iu\pi}\displaystyle
     \int_{C_r}\dfrac{f(w)}{w-z}\,\dd w$};
\end{tikzpicture}

{\small 柯西公式：全纯函数在圆盘
\emph{内部}的值是其在边界圆周上取值的加权平均。
全纯性所蕴含的一切刚性性质——
解析性、刘维尔定理与最大模原理，皆由
这一个恒等式推出。}
\end{omfigure}

\begin{method}\label{met:b3:holomorphic:toolkit}
日常工具箱。证明一个函数全纯：将其表示为
幂级数、复合、局部一致极限
（\cref{thm:b3:holomorphic:weierstrassconv}），或含全纯参数的积分
（见\cref{exo:b3:holomorphic:7}——
在 $\int$ 号下求导，或用莫雷拉–富比尼方法）。证明恒等式：
在线段或子区域上证明，再调用
恒等定理。估计界时，用柯西估计，取
可用的最大圆周。证明常值性或不存在性：
用刘维尔定理或最大模原理。始终要清楚你的函数在\emph{何处}
全纯，以及\emph{哪些}圆盘能放入
$\Omega$ 之中。
\end{method}

\section{练习}

\begin{exercise}[$\star$]\label{exo:b3:holomorphic:1}
(a) 在哪些点上 $z \mapsto \bar z$、$\abs z^2$、
$\operatorname{Re}z$ 在何处复可微？
在开集上是否全纯？
(b) 证明 $P(x, y) = x^2 - y^2$ 是某个全纯函数的实部，
该全纯函数定义于 $\C$ 上，请显式求出，
并确定\emph{全部}这样的函数。
\end{exercise}

\begin{exercise}[$\star$]\label{exo:b3:holomorphic:2}
由定义计算 $\int_{C}z^n\,\dd z$，对所有 $n
\in \Z$，$C$ 为单位圆周；$\int_\gamma\bar z\,\dd z$ 沿
线段 $[0, 1+\iu]$ 及经过
$1$ 的两段折线路径上的积分：由此得出 $\bar z$ 在
这些路径的任一邻域上都没有原函数。
\end{exercise}

\begin{exercise}[$\star\star$]\label{exo:b3:holomorphic:3}
(a) 证明主值对数 $\log z = \ln\abs z +
\iu\arg z$（$\arg \in \intoo{-\pi}\pi$）在
$\C\setminus\intoc{-\infty}0$ 上全纯，导数为 $\frac1z$
\emph{（在星形割平面上 $\frac1z$ 的原函数：
\cref{thm:b3:holomorphic:cauchy}；定出常数）}。
(b) 证明在
$\C^*$ 上不存在连续的对数 \emph{（其障碍与导数无关：
单位圆周的环绕数）}。
(c) 将 $\log(1 + z)$ 在 $D(0,1)$ 上展开为幂级数。
\end{exercise}

\begin{exercise}[$\star\star$]\label{exo:b3:holomorphic:4}
(a) 设 $f$ 为整函数，且 $\abs{f(z)} \leq C(1 + \abs
z)^{n}$。证明 $f$ 为次数 $\leq n$ 的多项式
\emph{（在大圆周上用柯西估计）}。
(b) 设 $f$ 为整函数，且 $\operatorname{Re}f$ 有上界。
证明 $f$ 为常值 \emph{（考虑 $\eu^{f}$）}。
(c) 由此推出“仿射情形的小皮卡定理”：整函数
若不取某一半平面内的值，则为常值。
\end{exercise}

\begin{exercise}[$\star\star$]\label{exo:b3:holomorphic:5}
(a) 设 $f$ 为连通开集 $\Omega \ni 0$ 上的全纯函数，且
$f(\frac1n) = \frac1{n^2}$ 对所有充分大的 $n$ 成立。求 $f$。
(b) 是否存在全纯函数 $f$（定义在 $\C^*$ 上），使得 $f(\frac1n) =
\frac{(-1)^n}{n}$ 对所有 $n \geq 1$ 成立？请说明理由。
(c) 给出两个不同的全纯函数，它们定义在
$D(0,1)\sqcup D(3,1)$ 上，且在 $D(0,1)$ 上相等：
恒等定理在何处用到了连通性？
\end{exercise}

\begin{exercise}[$\star\star$]\label{exo:b3:holomorphic:6}
设 $f$ 在开单位圆盘 $\mathbb D$ 上全纯，
在 $\bar{\mathbb D}$ 上连续，并且在
边界圆周上 $\abs f \equiv 1$。
(a) 若 $f$ 在 $\mathbb D$ 内无零点，证明 $f$ 为常值
\emph{（对 $f$ 与 $1/f$ 分别应用最大模原理）}。
(b) 举一个有零点但 $f$ 不为常值的例子。
\end{exercise}

\begin{exercise}[$\star\star$]\label{exo:b3:holomorphic:7}
（积分号下的全纯性）设 $\mu$ 为有限测度
在空间 $X$ 上，$g \colon X\times\Omega \to \C$ 满足：
$g(x, \cdot) \in \mathcal H(\Omega)$（对每个 $x$），$g$
关于 $x$ 可测，且 $\abs g \leq h(x)$，其中 $h$ 可积，
并且关于 $z$ 局部一致。证明 $G(z) = \int_Xg(x, z)\dd\mu(x)$
在 $\Omega$ 上全纯。\emph{（莫雷拉方法：三角形上的积分
由富比尼定理与古尔萨定理为零；连续性由控制收敛定理得出。
然后应用于 $\Gamma(z) =
\int_0^\infty t^{z-1}\eu^{-t}\dd t$ 在 $\{\operatorname{Re}z >
0\}$ 上的伽马函数）}。
\end{exercise}

\begin{exercise}[$\star\star\star$]\label{exo:b3:holomorphic:8}
（高斯–卢卡斯定理）设 $P \in \C[X]$ 为非常值多项式。证明
$P'$ 的每个根都位于
$P$ 的根的凸包中。\emph{（写出 $\frac{P'}{P} = \sum_k\frac{m_k}{z - a_k}$，
其中 $z$ 是 $P'$ 的根，且不是 $P$ 的根，取共轭，
并读出一个凸组合。）}以
$P = z^3 - 1$ 为例说明。
\end{exercise}

\begin{exercise}[$\star\star\star$]\label{exo:b3:holomorphic:9}
设 $f$ 为整函数，且为双周期函数：$f(z + 1) = f(z + \iu)
= f(z)$ 对所有 $z$ 成立。证明 $f$ 为常值。\emph{（先在紧的基本正方形上估计
$f$，再推广到全平面；
用刘维尔定理。）}由此可知：非常值椭圆函数必有极点——
这正是\cref{ch:b3:residues}的主题。
\end{exercise}

\begin{exercise}[$\star\star$]\label{exo:b3:holomorphic:10}
(a) 设 $P = \operatorname{Re}f$，其中 $f$ 为全纯函数。证明它
满足平均值性质 $P(a) =
\frac1{2\pi}\int_0^{2\pi}P(a + r\eu^{\iu t})\dd t$，并且
\emph{调和}，即 $\partial^2_{xx}P + \partial^2_{yy}P = 0$
（对柯西–黎曼方程求导；
并用\cref{thm:b3:holomorphic:analytic}说明所需的光滑性）。
(b) 由此推出：在有界区域上，全纯函数实部的
最大值原理。
\end{exercise}

\begin{exercise}[$\star\star\star$]\label{exo:b3:holomorphic:11}
（施瓦茨反射）设 $\Omega^+ = \{z : \abs z < 1,\
\operatorname{Im}z > 0\}$，$I = \intoo{-1}1$，且 $f$
在 $\Omega^+$ 上全纯，在 $\Omega^+\cup I$ 上连续，
\emph{在 $I$ 上取实值}。定义
\[
F(z) = \begin{cases} f(z) & z \in \Omega^+\cup I,\\
\overline{f(\bar z)} & \bar z \in \Omega^+ . \end{cases}
\]
(a) 证明 $F$ 是良定义的，并在
$\Omega = \Omega^+\cup I\cup\Omega^-$ 上连续，且在
$\Omega^\pm$ 上全纯 \emph{（对 $\Omega^-$：验证柯西–黎曼方程
对 $\overline{f(\bar z)}$ 成立，或把 $f$ 展开为局部幂级数
再取共轭系数）}。
(b) 用莫雷拉判别法证明 $F$ 在整个 $\Omega$ 上全纯：
只需 $\int_{\partial T}F = 0$ 对每个
三角形 $T \subseteq \Omega$ 成立 \emph{（在 $I$ 处切分三角形，
并把其水平边推离实轴
距离 $\varepsilon$，并利用一致连续性）}。
(c) 由此推出：圆盘上的全纯函数若在某一直径上取实值，
则满足 $f(\bar z) = \overline{f(z)}$；另有
非常值全纯函数不可能在
（连通）定义域的任一非空开子集上取实值。
\end{exercise}

\begin{exercise}[$\star\star$]\label{exo:b3:holomorphic:12}
（复数毕达哥拉斯方程）求所有
\emph{整函数}对，满足 $f^2 + g^2 = 1$。
(a) 证明 $h = f + \iu g$ 为无零点的整函数，并
证明每个无零点整函数都可写成 $\eu^{\varphi}$，其中
$\varphi$ 为某个整函数 \emph{（$h'/h$ 为整函数，因而在星形区域
$\C$ 上有原函数；适当调整常数并
证明 $h\eu^{-\varphi}$ 为常数）}。
(b) 由此得 $f = \cos\varphi$，$g = \sin\varphi$，其中
$\varphi$ 为整函数，并验证其逆命题。
求 $f^2 + g^2 = 0$ 的所有整函数解。
\end{exercise}

\section{问题：两次证明代数基本定理}

\begin{problem}[{周末问题 --- $\C$ 是代数闭的：
达朗贝尔的证明、刘维尔的证明以及
收获}]\label{pb:b3:holomorphic:1}
设 $P(z) = z^n + a_{n-1}z^{n-1} + \dots + a_0$，$n \geq 1$。
下面两次证明 $P$ 有根，然后汇集代数学
自\cref{ch:b3:galois}以来一直在等待的结论。

\medskip
\noindent\textbf{第 I 部分 --- 强制性与最小值。}
\begin{enumerate}
  \item 证明 $\abs{P(z)} \to +\infty$，当 $\abs z \to
        \infty$，精确地说，$\abs{P(z)} \geq \frac12\abs z^n$
        对适当选取的 $\abs z \geq R_0$ 成立。
  \item 由此可得 $\abs P$ 在 $\C$ 上取得全局最小值：
        即存在 $z_0$，使 $\abs{P(z_0)} = \inf_\C\abs P$
        \emph{（大闭圆盘的紧性，
        见\cref{cor:b3:topology:heineborel}）}。
\end{enumerate}

\medskip
\noindent\textbf{第 II 部分 --- 达朗贝尔下降法。}假设
用反证法，$P(z_0) \neq 0$。
\begin{enumerate}[resume]
  \item 将 $Q(h) = P(z_0 + h)/P(z_0)$ 展开为
        $h$ 的多项式：$Q(h) = 1 + c_kh^k + h^{k+1}S(h)$，其中 $c_k \neq
        0$，$k \geq 1$，$S$ 为多项式。
  \item 选取下降方向：对充分小的 $t > 0$，
        令 $h = t\,\omega$，其中 $\omega^k = -1/c_k$（为何存在这样的 $\omega$？
        --- \emph{证明}任一复数的
        $k$ 次方根存在，可借助极坐标形式，
        且不依赖于正在证明的定理
        ）。证明
        \[
        \abs{Q(t\omega)} \leq 1 - t^k + C\,t^{k+1}
        \]
        对充分小的 $t$，$C$ 为显式常数。
  \item 由此得 $\abs{Q(t\omega)} < 1$，对充分小的 $t$ ---
        与 $\abs{P(z_0)}$ 的极小性相矛盾。因此
        $P(z_0) = 0$：\emph{每个非常数复系数多项式
        都有根}（达朗贝尔--阿尔冈）。
\end{enumerate}

\medskip
\noindent\textbf{第 III 部分 --- 刘维尔的“一行证明”，
完整版。}
\begin{enumerate}[resume]
  \item 仔细写出
        \cref{cor:b3:holomorphic:liouville}的证明：若 $P$ 无根，
        验证 $1/P$ 为整函数，且有界（需给出定量估计，
        用到第 1 问），从而为常数，并由此得出结论。
        比较两种证明：每种证明各用到哪些要素？
        （紧性在两种证明中都出现 --- 出现在何处？）
\end{enumerate}

\medskip
\noindent\textbf{第 IV 部分 --- 收获。}
\begin{enumerate}[resume]
  \item 证明每个 $P \in \C[X]$（次数为 $n$）都可分解：
        $P = c\prod_{i}(X - \alpha_i)^{m_i}$，其中 $\sum m_i =
        n$，\emph{（归纳法，用 $(X -
        \alpha)$ 作带余除法）}。
  \item 证明 $\R[X]$ 中的不可约多项式恰为
        一次多项式，以及判别式为负的二次
        多项式 \emph{（共轭根成对）}；由此推出
        每个奇数次实多项式都有实根，
        并给出该结论的第二个证明，基于序结构
        （介值定理） ---
        并在 $X^3 - X - 1$ 上核对两者结论一致。
  \item 由此偿还本书现在能够偿还的“欠账”：(i) 每个
        非零有限维
        $\C$ 向量空间上的自同态都有特征值，因而每个复矩阵
        都有若尔当标准形
        （\cref{thm:b3:modules:jordan}）；(ii) 代数数域
        $\bar\Q$ 在
        \cref{rem:b3:galois:closureexamples}中所用，确实是
        $\Q$ 的代数闭包。
  \item （终章）指出各证明在一个
        类似 $\Q(\iu)$ 的域上会在何处失效：哪些步骤用到了
        $k$ 次方根的存在性（第 4 问），哪些用到了紧性
        或完备性（第 2 问与第 6 问）？用五行给出
        结论：该定理本质上是\emph{解析的} ---
        每个证明都在某处用到了实数的完备性或
        $\R$ 的连通性 --- 尽管其陈述是
        纯代数的。
\end{enumerate}

\medskip
\noindent\textbf{第 V 部分 --- 整函数的
刚性阶梯。}刘维尔定理是这一阶梯的第一级，我们将攀登
它。
\begin{enumerate}[resume]
  \item （柯西估计）由
        半径为 $r$、圆心为 $a$ 的圆周上的柯西公式，证明
        \[
        \bigl|f^{(n)}(a)\bigr| \leq
        \frac{n!\,\sup_{\abs{z-a}=r}\abs f}{r^n} ,
        \]
        并在 $n = 1$、$r \to
        \infty$ 的情形下恢复刘维尔定理。
  \item （多项式增长）证明：若整函数 $f$
        满足 $\abs{f(z)} \leq A + B\abs z^m$（对所有 $z$），则
        是次数 $\leq m$ 的多项式 \emph{（去掉 $m$ 次以上的泰勒系数，
        用第 11 问）}。
  \item （实部有界）证明：若整函数 $f$
        满足 $\operatorname{Re}f$ 有上界，则它为常数
        \emph{（对 $\eu^{f}$ 应用刘维尔定理）}。
  \item （双周期性）设整函数 $f$ 满足 $f(z + 1)
        = f(z)$，且 $f(z + \iu) = f(z)$ 对所有 $z$ 成立。证明 $f$
        为常数。由此可得：非常数的“椭圆”
        函数必有奇点 --- 历史上，
        极点正是由此进入复分析的。
  \item （稠密值域）证明非常数整函数的值域
        在 $\C$ 中稠密：若 $f(\C)$ 避开某圆盘
        $D(a, r)$，则 $\frac1{f - a}$ 为整函数且
        有界。（皮卡证明值域至多漏掉
        \emph{一个点}；稠密性是我们工具所能
        达到的水平。）
  \item （真性 $\Rightarrow$ 多项式）设 $f$ 为
        整函数，且 $\abs{f(z)} \to \infty$（当 $\abs z \to
        \infty$ 时）。证明：$f$ 的零点个数有限
        （$z_1, \dots, z_p$，重数分别为 $m_i$）；
        商 $g = f/\prod(z - z_i)^{m_i}$ 为整函数，且
        无零点；$1/g$ 具有多项式增长，因此
        （由第 12 问）为多项式，且必为常数
        （因其无零点）；由此得 $f$ 为多项式。于是
        在整函数中，多项式恰好就是
        真性的那些函数 --- $\eu^z$ 在
        $\R_-$ 上不满足真性。
\end{enumerate}

\medskip
\noindent\textbf{第六部分 --- 调和影子与
高斯的一个均值。}
\begin{enumerate}[resume]
  \item 设 $f = u + \iu v$ 在一个开集上全纯。
        验证 $u = \operatorname{Re}f$ 满足
        均值性质
        \[
        u(a) = \frac1{2\pi}\int_0^{2\pi}
        u\bigl(a + r\eu^{\iu\theta}\bigr)\,\dd\theta
        \]
        （柯西公式的实部），并由此推出
        有界区域上 $u$ 的最大值原理，
        其连通性证明与 $\abs f$ 的情形相同。
  \item （高斯均值）设 $a \in \C$，$r > 0$，且
        $\abs a \neq r$，证明
        \[
        \frac1{2\pi}\int_0^{2\pi}
        \log\bigl|a - r\eu^{\iu\theta}\bigr|\,\dd\theta
        = \log\max\bigl(\abs a, r\bigr)
        \]
        \emph{（若 $\abs a > r$：$z \mapsto \log\abs{a - z}$
        是闭圆盘某邻域上一个全纯对数的实部
        ——它为何存在？——
        因此可用第 17 问；若 $\abs a < r$：
        作变形 $\abs{a - r\eu^{\iu\theta}} = r\,\abs{1 -
        \frac ar\eu^{-\iu\theta}}$，并沿用前一种情形
        的论证）}。
  \item （马勒测度）设 $P = c\prod_{i=1}^n(X -
        \alpha_i) \in \C[X]$，由此推出\emph{多项式的
        延森公式}：
        \[
        \frac1{2\pi}\int_0^{2\pi}\log\bigl|P(\eu^{\iu\theta})
        \bigr|\,\dd\theta
        = \log\Bigl(\abs c\prod_{i}\max(1,
        \abs{\alpha_i})\Bigr) :
        \]
        单位圆周上 $\abs P$ 的几何平均，
        读出了圆盘外的根。在 $P =
        X^2 - X$ 与 $P = 2X - 1$ 上验证。
  \item （伯努利数）定义系数 $B_n$，使得
        $\frac{z}{\eu^z - 1} = \sum_{n\geq0}\frac{B_n}{n!}
        z^n$ 在 $0$ 附近成立（为何左端在
        $0$ 处解析？）。导出递推式
        $\sum_{k=0}^{n}\binom{n+1}kB_k = 0$（$n \geq 1$），依据是
        $(\eu^z - 1)\cdot\frac z{\eu^z-1} = z$，计算 $B_0,
        \dots, B_6$，并证明 $B_{2k+1} = 0$（$k \geq 1$）
        \emph{（函数 $\frac z{\eu^z-1} + \frac z2$ 是
        偶函数）}。这些数将给出
        在\cref{ch:b3:residues}中每个 $\zeta(2k)$ 的值。
  \item （实性）证明：在 $\R$ 上取实值的整函数
        满足 $f(\bar z) =
        \overline{f(z)}$ 在全平面上成立 \emph{（比较
        $0$ 处的泰勒系数，或对
        $z \mapsto \overline{f(\bar z)}$ 应用唯一性定理）}；再次
        推出实系数多项式的非实根成共轭对出现
        （第 8 问的配对，此处以分析方法再证
        一遍）。
  \item （总结）梳理刚性阶梯：有界
        $\Rightarrow$ 常数；多项式有界
        $\Rightarrow$ 多项式；真性 $\Rightarrow$
        多项式；值域缺少一个圆盘 $\Rightarrow$ 常数；
        双周期 $\Rightarrow$ 常数。以一小段文字，将其与实
        $\mathcal C^\infty$ 函数作对比
        （隆起函数，
        \cref{thm:b3:lp:regularization}）：为何
        全纯性这一纯局部条件，会施加全局的
        规律与秩序？
\end{enumerate}

\medskip
\noindent\textbf{第七部分 --- 最后的收获。}
\begin{enumerate}[resume]
  \item （兰道不等式）设 $P = \sum_{k=0}^na_kX^k$，
        证明均值 $\frac1{2\pi}\int_0^{2\pi}
        \abs{P(\eu^{\iu\theta})}^2\dd\theta =
        \sum_k\abs{a_k}^2$（借助
        $\eu^{\iu k\theta}$ 的正交性），再用逐点估计
        $\log t \leq t - 1$ 来比较
        $\log\abs P^2$ 与 $\abs P^2$ 的均值，由第
        19 问可得：
        \[
        \abs c\prod_{i}\max\bigl(1, \abs{\alpha_i}\bigr)
        \;\leq\;
        \Bigl(\sum_{k=0}^{n}\abs{a_k}^2\Bigr)^{1/2} :
        \]
        单位圆盘外各根之积
        受系数的 $\ell^2$ 大小控制
        （位于圆周\emph{上}的根，可对
        $P(rX)$ 应用该不等式并令 $r \to 1$）。
        在 $X^2 - X$ 上验证此式。
  \item （伯努利数按阶乘增长）证明
        $\sum\frac{B_n}{n!}z^n$ 的收敛半径
        \emph{恰为} $2\pi$：至少为 $2\pi$，因为
        $z/(\eu^z - 1)$ 可全纯延拓至
        $D(0, 2\pi)$，至多为 $2\pi$，否则级数和
        在 $2\pi\iu$ 附近将保持有界，而那里
        $\abs{z/(\eu^z-1)} \to \infty$，由此可得
        \[
        \limsup_{k\to\infty}
        \Bigl(\frac{\abs{B_{2k}}}{(2k)!}\Bigr)^{1/2k}
        = \frac1{2\pi} :
        \]
        伯努利数按阶乘增长。承认 $B_{12}
        = -\frac{691}{2730}$（即第 20 问的递推式
        继续算下去），比较 $\abs{B_{12}} \approx 0.25311$
        与更精细的预测 $2\,(2k)!/(2\pi)^{2k}
        \approx 0.25305$（$k = 6$ 时）——四位数字吻合，
        这一渐近律，\cref{ch:b3:residues} 将
        借助 $\zeta(2k)$ 精确证明之。
  \item （根的连续移动）设 $(P_j)$ 为
        $n$ 次首一多项式，其系数收敛于
        （首一）$P$ 的系数。证明\emph{柯西界}：
        首一多项式 $Q = X^n + \sum_{k<n}q_kX^k$ 的每个根
        满足 $\abs\alpha \leq 1 + \max_k\abs{q_k}$；由此推出
        $P_j$ 的根落在一个固定的
        紧集内，并抽取收敛的
        根向量子列，并在第
        7 问的因式分解中取极限，得出
        $P_j$ 的根重集收敛于 $P$ 的根重集。
        最后证明，连续性是人们所能得到的最佳
        结论：例如对 $P_\varepsilon = X^2 - 2X + 1 +
        \varepsilon$，一个大小为 $\varepsilon =
        10^{-4}$ 的扰动使二重根 $1$ 移动 $10^{-2}$ ——
        赫尔德指数为 $\frac1m$，位于 $m$ 重根处，
        绝不是利普希茨的：数值上看，重根
        损失一半的有效数字。
\end{enumerate}
\end{problem}
